\documentclass{article}
\usepackage[utf8]{inputenc}
\usepackage{fullpage}
\setcounter{tocdepth}{2}
\usepackage{hyperref}
\usepackage{graphicx}
\usepackage{float}

% \usepackage{apalike}
% \usepackage{natbib}
\usepackage[style=apa]{biblatex}
\addbibresource{references.bib}

\newcommand{\fcite}[1]{\footnote{\fullcite{#1}}}

\title{Kinship Structures} 

\author{EDSD Course syllabus}

\date{Last updated: \today}

\begin{document}

\maketitle

\noindent
European Doctoral School of Demography (EDSD)\\
Paris, March 2-5, 2026\\
% Lectures: from Monday to Thursday, 10:00 – 11:30\\
% Lectures: from Monday to Thursday, 10:00 – 11:30\\

\noindent
Instructors:\\
Diego Alburez-Gutierrez\\
Kinship Inequalities Research Group\\
Max Planck Institute for Demographic Research\\
\href{mailto:alburezgutierrez@demogr.mpg.de}{alburezgutierrez@demogr.mpg.de}\\


% \tableofcontents

\section{Introduction}

\subsection{Course description}

Kinship is a fundamental property of human populations and a key form of social structure.
Demographers have long been interested in the interplay between demographic change and family configuration.
This has led to the development of sophisticated methodological and conceptual approaches for the study of kinship, some of which are reviewed in this course. 

\subsection{Goals}

The course will provide a succinct and practical introduction to kinship demography. 
By the end of the course, students will be familiar with the current state of the sub-field and the tools to conduct independent research in this area. 
The main goals of the course are to:

\begin{enumerate}
    \item present a general overview of the current state of the field of kinship demography,
    \item introduce the fundamentals of the formal demography of kinship, and
    \item provide hands-on experience running models of kinship in R.
\end{enumerate}

\subsection{Lectures and exercises}

The in-person course comprises independent reading, lectures, and lab sessions. There are required and optional readings each day. Students are expected to attend all lectures and lab sessions. There is no exam; only a final assignment which needs to be turned in at the end of the week (see Section \ref{sec:assignment}).  

\subsection*{Schedule}

\begin{itemize}
    \item In-person sessions take place between 09:00 and 12:30 (note that on Monday we will start at 09:30 with the course content)
    \item From Monday to Wednesday, we will have 1.5 hours of lectures and 1.5 hours of lab session (with breaks in between)
    \item On Thursday, we will meet to work on your assignments 
\end{itemize}

\subsection{Hardware and software}

Lab sessions will be in R ($\geq4.0.2$). 
Students are required to bring a laptop for the lab sessions. 
Participants should install the following packages from CRAN: \texttt{DemoKin}, \texttt{dplyr}, \texttt{tidyr}, \texttt{ggplot2}, \texttt{readr}, and \texttt{data.table}.

\subsection{Online resources and lab sessions}

There are two online resources with all the information you need about the course, including lab sessions. We will update the materials constantly, so please make sure that you are working with the latest version of the data.:

\begin{itemize}
    \item A GitHub repository including the syllabus, slides, and other important materials: \url{https://github.com/alburezg/EDSD_2025_kinship}. Please make sure to download (clone) the \textbf{entire repository} into your computer!
    \item The website with the lab exercises: \url{https://alburezg.github.io/EDSD_2025_kinship/}.
\end{itemize}

\section{Final assignment}
\label{sec:assignment}

The assignment should be completed in groups (of three) that will be defined at the start of the course.

\subsection{Description}

You will use data on kinship structures to benchmark formal models of kinship.
For this exercise, you will use the \texttt{DemoKin} R package to implement formal models of kinship. 
You should choose one country and run four different models according to the following specifications:

    \begin{itemize}
        \item One-sex model; approximate male kin using GKP factors
       \begin{itemize}
           \item time-invariant rates
           \item time-variant rates
       \end{itemize}  
        
        \item Two-sex model; approximate male kin using the androgynous assumption
         \begin{itemize}
           \item time-invariant rates
           \item time-variant rates
             \end{itemize}  
     \end{itemize}

    Use the output of the four models to answer the following questions:

    \begin{enumerate}
     
    \item Plot the \textbf{expected number of living relatives by age of focal} for each specification. For extra points (i.e., this is optional), also plot the \textbf{expected number of deceased relatives by age of focal}.
    \item Discuss 1-2 key insights, when would you use different specifications? Consider the specific context and the data available for the country you selected. (max 250 words)
    \item Can you think of other ways of incorporating male fertility into the kinship models (beyond the options we discussed in the course)?  (max 250 words)
       \end{enumerate}  


% \subsection{Data}

% \begin{itemize}
%     \item We can obtain rates from any country in the world from United Nations 2024 Revision of the World Population Prospects. Pick any country (feel free to pick at random or pick the one you are interested in) and download the following data:
%     \begin{enumerate}
%         \item \href{https://population.un.org/wpp/Download/Files/1_Indicators%20(Standard)/CSV_FILES/WPP2022_Fertility_by_Age1.zip}{Female age-specific fertility rates, 1950 - 2100}
%         \item \href{https://population.un.org/wpp/Download/Files/1_Indicators%20(Standard)/CSV_FILES/WPP2022_Life_Table_Complete_Medium_Female_1950-2021.zip}{Female Life Tables, 1950 - 2021}
%         \item \href{https://population.un.org/wpp/Download/Files/1_Indicators%20(Standard)/CSV_FILES/WPP2022_Life_Table_Complete_Medium_Male_1950-2021.zip}{Male Life Tables, 1950 - 2021}        
%     \end{enumerate}
%         \end{itemize}

\subsection{Handing in the assignment}

Assignments (one per group) should be sent to \href{mailto:alburezgutierrez@demogr.mpg.de}{alburezgutierrez@demogr.mpg.de} \textbf{before midnight of Friday, March 6.} 
You should hand in the following files:
\begin{enumerate}
    \item An .RMD file with all your code and answers to the exercise questions
    \item A compiled .pdf of your markdown file showing all the code
    \item All input data needed to replicate your code 
\end{enumerate}


\section{Lecture Plan}

\subsection*{Monday, Mar 2: Introduction to kinship demography}

\subsubsection*{Required reading}

\begin{itemize}
    \item \fullcite{alburez-gutierrez_kinship_2022}
\end{itemize}

% \subsubsection*{Optional reading}

% \begin{itemize}
%     \item \fullcite{brass_derivation_1953}
%     \item \fullcite{bongaarts_projection_1987}
%     \item \fullcite{alburez-gutierrez_womens_2021}
% \end{itemize}

\subsection*{Tuesday, Mar 3: Mathematical models of kinship}

\subsubsection*{Required reading}

\begin{itemize}
    \item \fullcite{caswell_formal_2019}
\end{itemize}

\subsubsection*{Optional reading}

\begin{itemize}
    \item \fullcite{goodman_family_1974}
    % \item\fullcite{keyfitz_demographic_2005}    
\end{itemize}


\subsection*{Wednesday, Mar 4: Extensions of the kinship model}

\subsubsection*{Required reading}

\begin{itemize}
   \item \fullcite{alburez-gutierrez_projections_2023}
\end{itemize}

\subsubsection*{Optional reading}

\begin{itemize}
     \item \fullcite{caswell_formal_2021}
     \item \fullcite{caswell_formal_2022}
    % \item \fullcite{song_role_2022}    
\end{itemize}

\subsection*{Thursday, Mar 5: Meet to work on assignment}

\subsubsection*{No required readings}

\end{document}
